\section{Practical Test Expectations}

This section explains how the practical-test-related tutorial submissions are used in the module. You must submit all assigned weekly tutorial exercises through the relevant ITEL tabs, even though only selected activities are graded formally. These submissions help the lecturer verify your individual contribution and your understanding of the methods introduced in class.

The tutorial sessions that receive particular assessment attention are:
\begin{itemize}
    \item \textbf{Identifying Customer Needs};
    \item \textbf{Product Specifications}; and
    \item \textbf{Concept Generation}.
\end{itemize}

Use the official template provided on ITEL unless the lecturer announces a different format. A well-prepared submission should be clearly organised, properly captioned, and supported by appropriate references or evidence.

\subsection{What the Lecturer Will Look For}

The criteria below show how the submitted work supports the relevant \textbf{Course Learning Outcomes (CLOs)}.

\subsubsection{Identifying Customer Needs}
\begin{itemize}
    \item clear survey or interview objectives;
    \item relevant prompts and questions;
    \item accurate translation of customer statements into interpreted needs;
    \item sufficient breadth of needs and themes; and
    \item clear evidence that customer statements are authentic and relevant.
\end{itemize}

\subsubsection{Product Specifications}
\begin{itemize}
    \item a justified importance-versus-needs table;
    \item measurable, functionality-focused metrics;
    \item clear justification for selected units and metric relevance;
    \item a needs-to-metric matrix;
    \item a competitive benchmarking comparison supported by credible sources; and
    \item properly captioned tables and figures.
\end{itemize}

\subsubsection{Concept Generation}
\begin{itemize}
    \item clear analysis of primary functions;
    \item justification for the selected decomposition method;
    \item correct use of action-oriented functional descriptions;
    \item breadth and quality of solution options; and
    \item appropriate citation of internal and external sources.
\end{itemize}

\subsection{Reflection Prompt}

After each practical-test-related submission, briefly ask yourself: Which part of the method do I understand well, and which part still needs stronger evidence, better structure, or clearer explanation?
